\documentclass[11pt]{article}
\usepackage[T1]{fontenc}
\usepackage{textcomp}
\usepackage[margin=0.75in]{geometry}
\usepackage{amsmath,amssymb,amsthm}
\usepackage{array,booktabs}
\usepackage{hyperref}
\setlength{\parindent}{0pt}
\newtheorem{theorem}{Theorem}
\theoremstyle{definition}
\newtheorem{definition}{Definition}
\title{Flow Proof Artifact: List-rev-rev}
\author{Generated by \texttt{flow doc proof}}
\date{Flow Proof Book}
\begin{document}
\maketitle
Reversing twice recovers the original list.\\[0.5em]
\textbf{Source.} church — https://en.wikipedia.org/wiki/List\_(abstract\_data\_type)\\[0.5em]
\bigskip
\noindent{\large \textbf{Derived fact 1.} \textit{rev(rev(xs)) = xs} \hfill List \cdot reverse \cdot double reverse returns the list}\par
\smallskip\noindent\textit{Coordinate: List \cdot reverse \cdot double reverse returns the list}\par
\smallskip\noindent\textit{Source: church/induction}\par
\smallskip\noindent\textit{Built on: reversing empty yields empty, for reverse on List, reverse distributes over append, for reverse on List}\par
\medskip
\noindent\textbf{Goal.} rev(rev(xs)) = xs\par
\noindent$\displaystyle \forall xs \in List\quad rev(rev(xs)) = xs$\par
\medskip
\noindent
\begin{tabular}{@{} >{\raggedright\arraybackslash}p{0.06\textwidth} >{\raggedright\arraybackslash}p{0.47\textwidth} >{\raggedleft\arraybackslash}p{0.41\textwidth} @{}}
\textbf{\#} & \textbf{Proof} & \textbf{Mathematics} \\
\midrule
\textbf{1.} & We prove that double reverse returns the list for reverse on List. & \\[0.45em]
\textbf{2.} & We proceed by induction on xs: first the base case, then the inductive step. & \\[0.45em]
\textbf{3.} & Consider the base case in which xs  equals  nil. & \\[0.45em]
\textbf{4.} & We invoke the derived fact governing reverse on List: reversing empty yields empty, for reverse on List. & \\[0.45em]
\textbf{5.} & From step 3 and step 4, we can deduce that rev(rev(xs)) equals xs. This establishes the base case (see step 3 and step 4). Hence proven. & $\displaystyle rev(rev(xs)) = xs$ \\[0.45em]
\textbf{6.} & For the inductive step, suppose the claim holds for all smaller values. & \\[0.45em]
\textbf{7.} & Under the supposition in step 6, let rest = tail(xs). & \\[0.45em]
\textbf{8.} & Under the supposition in step 6, we cross the inductive boundary: assume the claim holds for rest (the induction hypothesis). & $\displaystyle rev(rev(rest)) = rest$ \\[0.45em]
\textbf{9.} & We invoke the derived fact governing reverse on List: reverse distributes over append, for reverse on List (instantiated for rest, cons(head(xs), nil)). & \\[0.45em]
\textbf{10.} & From step 6, step 7, step 8, and step 9, this implies rev(rev(xs)) equals xs. Hence proven. & $\displaystyle rev(rev(xs)) = xs$ \\[0.45em]
\bottomrule
\end{tabular}
\label{thm:List-reverse-double_reverse_returns_the_list}
\par\medskip\hrule
\end{document}
